\section*{第 \ref{ch:b2:biogeochemical-cycles} 章——生物地球化学循环}

\begin{solution}{exo:b2:biogeochemical-cycles:1}
\omterm{def:b2:biogeochemical-cycles:reservoir}{储库}是某种元素的储存量（质量）；\omterm{def:b2:biogeochemical-cycles:reservoir}{通量}是\omterm{def:b2:biogeochemical-cycles:reservoir}{储库}之间的
转移速率（质量每年）；\omterm{def:b2:biogeochemical-cycles:reservoir}{停留时间}是
\omterm{def:b2:biogeochemical-cycles:reservoir}{储库}的质量除以总输出量，即一个原子
在此停留的平均时间。陆地植物：$450/120 \approx 4$ 年（大致指叶片的
碳；\omterm{def:b2:plant-meristems:wood}{木材}可留存数十年，平均值掩盖了
分布差异）。深海：$37000/100 \approx 400$ 年——指
一个已沉降的分子再次见到表层之前所需的时间。
\end{solution}

\begin{solution}{exo:b2:biogeochemical-cycles:2}
$2.5\times 2.13 = \qty{5.3}{GtC}$ 每年。$10/2.13 = \qty{4.7}{ppm}$
——若全部留在大气中；按 \qty{45}{\%} 的大气留存比例，
则为 \qty{2.1}{ppm}。
\end{solution}

\begin{solution}{exo:b2:biogeochemical-cycles:3}
输入：\omterm{prop:b2:biogeochemical-cycles:nitrogen}{固氮}酶介导的生物\omterm{prop:b2:biogeochemical-cycles:nitrogen}{固氮}（蓝细菌、
根瘤菌、\emph{Azotobacter}、\emph{Frankia}）和闪电。返回：
\omterm{def:b2:microbial-metabolism:anaerobic}{反硝化作用}（兼性厌氧菌，如 \emph{Pseudomonas}
和 \emph{Paracoccus}，见于缺氧土壤和沉积物中）以及\omterm{prop:b2:biogeochemical-cycles:nitrogen}{厌氧氨氧化}
（浮霉菌）。中间过程：分解者（真菌、细菌）将有机氮
转化为铵（氨化作用）；铵经\omterm{prop:b2:microbial-metabolism:lithotrophy}{硝化作用}
氧化为亚硝酸盐（\emph{Nitrosomonas}、
氨氧化古菌），亚硝酸盐再氧化为硝酸盐（\emph{Nitrobacter}）；以及
植物和微生物对铵或硝酸盐的同化作用。
\end{solution}

\begin{solution}{exo:b2:biogeochemical-cycles:4}
在地质时间尺度上，氮总可由\omterm{prop:b2:biogeochemical-cycles:nitrogen}{固氮}生物
从空气中取之不尽的 $\mathrm{N_2}$ 制得，每当氮短缺时，\omterm{prop:b2:biogeochemical-cycles:nitrogen}{固氮}者便受青睐；
磷没有这样的\omterm{def:b2:biogeochemical-cycles:reservoir}{储库}，其供应
由风化作用决定，因此长期来看生产力的上限
是磷。一个季节之内，\omterm{prop:b2:biogeochemical-cycles:nitrogen}{固氮}缓慢且代价高，
多数植物无法\omterm{prop:b2:biogeochemical-cycles:nitrogen}{固氮}，因此眼下可得的氮
正是限制此时此地生长的因素——这也是化肥
以氮为主的原因。
\end{solution}

\begin{solution}{exo:b2:biogeochemical-cycles:5}
$\tau = 1/0.2 = 5$ 年；$M^{*} = F\tau = \qty{10}{t}$，即
$\qty{1}{g/m^3} = \qty{1000}{mg/m^3}$，严重\omterm{prop:b2:biogeochemical-cycles:ps}{富营养化}。达到
\qty{90}{\%}所需的时间为 $\tau\ln 10 = 11.5$ 年。输入减半会使
\omterm{def:b2:biogeochemical-cycles:reservoir}{稳态}减半至 \qty{500}{mg/m^3}，
并在相同的时间尺度上达到。
\end{solution}

\begin{solution}{exo:b2:biogeochemical-cycles:6}
1960 年代：\qty{0.9}{ppm/yr}，\qty{1.9}{GtC/yr}。2010 年代：\qty{2.3}{ppm/yr}，
\qty{5.0}{GtC/yr}。增长率提高了
两倍半，与排放量同步增长。
\end{solution}

\begin{solution}{exo:b2:biogeochemical-cycles:7}
\qty{6}{ppm} 相当于 \qty{13}{GtC}，而净初级生产力为
\qty{35}{GtC}：锯齿形波动为其三分之一。它之所以较小，是因为
呼吸作用与分解作用在整个夏季仍在进行（锯齿
记录的是净吸收而非总吸收），
南半球的季节相反、部分相互抵消，
以及海洋会抑制这一波动。
\end{solution}

\begin{solution}{exo:b2:biogeochemical-cycles:8}
$\qty{200}{kg}/(\qty{28}{g/mol}) = 7100$~摩尔的 $\mathrm{N_2}$；$16\times
7100 = \num{114000}$~摩尔 ATP；$/30 = 3800$~摩尔葡萄糖，\qty{690}{kg}。
一种积累 \qty{10}{t} 干物质的作物，计入呼吸作用后，
约固定 \qtyrange{15}{20}{t} 糖，因此\omterm{prop:b2:biogeochemical-cycles:nitrogen}{固氮}耗去其
光合产物的约 \qty{4}{\%}——这便是
摆脱土壤氮依赖的代价。
\end{solution}

\begin{solution}{exo:b2:biogeochemical-cycles:9}
可用比值 $\mathrm{N:P} = 32/2.2 = 14.5$，低于雷德菲尔德比
16：硝酸盐首先耗尽，磷酸盐尚余
\qty{0.2}{\micro mol/kg}。固定的碳：$32\times 106/16 =
\qty{212}{\micro mol/kg}$，约合每千克水含
碳 \qty{2.5}{mg}。
\end{solution}

\begin{solution}{exo:b2:biogeochemical-cycles:10}
$E^{*} = Q\tau = \qty{500}{GtC}$，比工业化前的 285
高出 \qty{235}{ppm}，约 \qty{520}{ppm}。实际情况更糟，因为
这些库不是单一的快速池：表层海洋趋于饱和
（其吸收 $\mathrm{CO_2}$ 的能力随碳酸盐被消耗而下降），
陆地库随变暖而减弱，而真正的移除过程还有
一个缓慢的组成部分——深海混合与岩石风化——其
时间常数为一千至十万年，因此过剩量的一部分
不会在任何几十年量级的 $\tau$ 内弛豫。
\end{solution}

\begin{solution}{exo:b2:biogeochemical-cycles:11}
所测量的是大气中 $\mathrm{CO_2}$ 与
海洋和生物圈之间的交换，而非衰变：核试验释放的 $\mathrm{{}^{14}C}$ 被稀释
进入大得多的表层海洋库及植物
和土壤库之中，此后它们又释放出普通碳。超量部分的 11 年半衰期
（e 折合时间约 16 年）即为
该交换的时间尺度——也就是大气碳相对于表层海洋与陆地的
周转时间，比
四年的总\omterm{def:b2:biogeochemical-cycles:reservoir}{停留时间}更长，因为离开大气的碳有很大一部分
会直接返回。
\end{solution}

\begin{solution}{exo:b2:biogeochemical-cycles:12}
氮：自然\omterm{prop:b2:biogeochemical-cycles:nitrogen}{固氮}每年约 200~Tg；人为\omterm{prop:b2:biogeochemical-cycles:nitrogen}{固氮}
（哈伯–博施 120、作物 60、燃烧 30）约 210：\omterm{def:b2:biogeochemical-cycles:reservoir}{通量}
翻了一倍，受影响的\omterm{def:b2:biogeochemical-cycles:reservoir}{储库}——土壤、河流、近岸
海域以及大气中的 $\mathrm{N_2O}$——体量小而周转快，因此变化
在数十年内即可显现，而 $\mathrm{N_2}$ \omterm{def:b2:biogeochemical-cycles:reservoir}{储库}未受影响。
碳：人为排放每年 11~GtC，只占 \qty{5}{\%} 的
自然总交换量（约 200），因此\omterm{def:b2:biogeochemical-cycles:reservoir}{通量}只是略有扰动；
但自然\omterm{def:b2:biogeochemical-cycles:reservoir}{通量}相互抵消，人为\omterm{def:b2:biogeochemical-cycles:reservoir}{通量}却不然，并且它已使
大气\omterm{def:b2:biogeochemical-cycles:reservoir}{储库}增加了 \qty{50}{\%}——这一微小
扰动在一个最终库较慢的\omterm{def:b2:biogeochemical-cycles:reservoir}{储库}上产生了很大的累积效应。
哪一个循环“受扰动更大”，取决于人们
计的是\omterm{def:b2:biogeochemical-cycles:reservoir}{通量}还是累积量。
\end{solution}

\begin{solution}{pb:b2:biogeochemical-cycles:1}
\textbf{1.} $285\times 2.13 = \qty{607}{GtC}$; $412\times 2.13 =
\qty{878}{GtC}$；增加 \qty{271}{GtC}。
\textbf{2.} $271/700 = 0.39$（近几十年的比例，
计入土地利用排放时约为 $0.45$）。
\textbf{3.} 海洋吸收约四分之一的排放量，陆地
（森林再生与 $\mathrm{CO_2}$ 施肥效应）约占三分之一；
其余 \qty{430}{GtC} 未留在大气中。
\textbf{4.} $175/38000 = \qty{0.5}{\%}$——对整个海洋而言可忽略不计；
但迄今为止，吸收的碳主要进入表层
（\qty{900}{GtC}），在那里意味着升高了十分之一甚至更多，
碳酸盐化学过程又使溶解 $\mathrm{CO_2}$ 的小幅升高转化为
氢离子浓度的较大升高：表层 pH 下降了 $0.1$，
相当于酸度升高 \qty{30}{\%}。
\textbf{5.} 令 $Q = Q_0 e^{t/T}$，试取 $E = Ae^{t/T}$：$A/T = Q_0 -
A/\tau$，故 $A = Q_0\tau T/(\tau + T)$，大气留存比
$\dot E/Q = \tau/(\tau + T)$ 为常数。排放量以每年
\qty{2}{\%} 的速率增长（$T = 50$），取 $\tau = 40$，得 $0.44$：
留存比保持不变，因为库与源同步增长。
\textbf{6.} 若 $Q$ 为常数，则 $E \to Q\tau$，$\dot E \to 0$，
随着库逐渐赶上，大气留存比趋于零——这是在
单箱模型中的结果；实际上库会趋于饱和，留存比
下降得更慢。
\textbf{7.} \qty{271}{GtC}.
\textbf{8.} $E^{*} = 10\times 100 = \qty{1000}{GtC}$: $285 +
1000/2.13 = \qty{755}{ppm}$.
\textbf{9.} $E(t) = 1000 + (271 - 1000)e^{-t/100}$；当 $t = 80$ 时，
$E = 1000 - 729\times 0.449 = \qty{672}{GtC}$: \qty{600}{ppm}.
\textbf{10.} $Q = 10(1 - t/50)$：特解 $E_p = at + b$
其中 $a = -20$，$b = 3000$（由 $a = 10 - b/\tau$ 及 $0 = -0.2 -
a/\tau$）；$E = 3000 - 20t + (271 - 3000)e^{-t/100}$。当 $t = 50$ 时：
$2000 - 2729\times 0.607 = \qty{344}{GtC}$，\qty{447}{ppm}。此后 $E$
自由衰减：当 $t = 80$ 时，$344\,e^{-0.3} = \qty{255}{GtC}$，
\qty{405}{ppm}。
\textbf{11.} \qty{600}{ppm} 对比 \qty{405}{ppm}：第二种
情景使大气到 2100 年时几乎回到今天的水平。
\textbf{12.} 单一的 $\tau$ 使库在
同一速率下继续作用于不断缩小的过剩量；实际上快速的库（
混合层、正在生长的森林）已经吸收了它们能吸收的部分，
剩余的移除主要依靠缓慢的深海混合与
风化作用，其 $\tau$ 为数世纪至数千年，因此净零排放之后的下降
要慢得多——浓度在数世纪内大致保持
不变，而非下降。
\textbf{13.} 作物吸收 \qty{90}{kg}；反硝化 \qty{37.5}{kg}；浸出
\qty{22.5}{kg} 的氮（每公顷每年）。
\textbf{14.} $\qty{22.5}{kg}/\qty{14}{g/mol} = 1600$~摩尔 N；碳 
$1600\times 106/16 = \num{10600}$~摩尔，\qty{128}{kg}。
\textbf{15.} \num{10600}~摩尔 $\mathrm{O_2}$，\qty{340}{kg}。
\textbf{16.} $\num{10600}/0.25 = \num{42000}$~\unit{m^3}：一公顷
的浸出每年可使四万立方米底层水失去氧气——而
密西西比河汇集了
一亿公顷的流域。
\textbf{17.} $37.5\times 0.01 = \qty{0.375}{kg}$ 的氮以
$\mathrm{N_2O}$ 形式存在，即 $0.375\times 44/28 = \qty{0.59}{kg}$ 的 $\mathrm{N_2O}$；
$\times 270 = \qty{160}{kg}$ 的 $\mathrm{CO_2}$ 当量。
\textbf{18.} $150\times 4 = \qty{600}{kg}$ 的 $\mathrm{CO_2}$：
制造过程的排放量是田间 $\mathrm{N_2O}$ 的四倍，两者合计
每公顷每年为四分之三吨。
\textbf{19.} $\tau = 10$ 年；$M^{*} = 5\times 10 = \qty{50}{t}$；
浓度 $50/(5\times 10^{7}) = \qty{1}{g/m^3} =
\qty{1000}{mg/m^3}$。
\textbf{20.} 超过阈值三十倍：严重\omterm{prop:b2:biogeochemical-cycles:ps}{富营养化}。
\textbf{21.} $1/\tau = 0.1 + 0.2 = 0.3$：$\tau = 3.3$ 年；$M^{*} =
5/0.3 = \qty{16.7}{t}$, \qty{330}{mg/m^3}.
\textbf{22.} $M^{*} = 1/0.3 = \qty{3.3}{t}$，\qty{67}{mg/m^3}——仍然
\omterm{prop:b2:biogeochemical-cycles:ps}{富营养化}；在 $\tau\ln 10 = 7.7$ 年后降至 \qty{10}{\%} 以内。
\textbf{23.} 沉积物并非单向的库：\omterm{prop:b2:biogeochemical-cycles:ps}{富营养化}期间储存在其中的磷
会在\omterm{prop:b2:biogeochemical-cycles:ps}{富营养化}本身造成的缺氧条件下释放回水体
（与铁结合的磷酸基团在铁被还原时溶解），
因此湖泊在外源输入切断后仍能自我供给数年——
即内部负荷，一个
\omterm{def:b2:biogeochemical-cycles:reservoir}{停留时间}很长的第二\omterm{def:b2:biogeochemical-cycles:reservoir}{储库}。
\textbf{24.} $\qty{5000}{kg}/\qty{31}{g/mol} = 1.6\times 10^{5}$~摩尔 P；
碳 $\times 106 = 1.7\times 10^{7}$~摩尔，即每年 \qty{205}{t} 碳，
折合湖面约 \qty{20}{g/m^2}（若湖深 \qty{5}{m}），
属于严重的水华。
\textbf{25.} 1850--2020 年间大气留存比为 $0.39$；2100 年时，
排放量维持在 \qty{10}{GtC/yr} 时为 \qty{600}{ppm}，
而排放量于 2070 年前降至零（一种乐观的
单一 $\tau$ 模型）时为 \qty{405}{ppm}；氮浸出量为每公顷
每年 \qty{22.5}{kg}。
\end{solution}
